\documentclass[../../../include/open-logic-section]{subfiles}

\begin{document}

\olfileid[id]{his}{set}{pathology}
\olsection{Patologi}

Namun, ternyata definisi \emph{limit} memungkinkan beberapa konstruksi yang
cukup ``patologis''.

Sekitar dasawarsa 1830-an, Bolzano menemukan suatu fungsi yang
\emph{kontinu di mana-mana}, tetapi \emph{tidak dapat didiferensialkan di
mana pun}. (Sayangnya, Bolzano tidak pernah menerbitkan hasil ini; para
matematikawan baru menjumpai gagasan tersebut pada 1872 berkat penemuan
independen Weierstrass atas gagasan yang sama.)\footnote{Sejarahnya
didokumentasikan dengan sangat menyeluruh dalam catatan-catatan kaki artikel
Wikipedia tentang
\href{http://en.wikipedia.org/wiki/Weierstrass_function}{fungsi
Weierstrass}.} Setidaknya, hasil ini cukup mengejutkan. Mudah menemukan
fungsi seperti $|x|$ yang kontinu di mana-mana tetapi tidak dapat
didiferensialkan pada satu titik tertentu. Namun, suatu fungsi yang kontinu
di mana-mana tetapi \emph{tidak dapat didiferensialkan di mana pun} sungguh
merupakan makhluk yang sangat berbeda. Bayangkan sejenak bagaimana Anda
mencoba menggambar fungsi seperti itu. Untuk menjamin kontinuitasnya, Anda
harus dapat menggambarnya tanpa pernah mengangkat pena dari halaman; tetapi
untuk menjamin bahwa ia tidak dapat didiferensialkan di mana pun, Anda harus
terus-menerus mengubah arah pena secara mendadak.

``Patologi'' lain menyusul. Pada 5 Januari 1874, Cantor menulis surat
kepada Dedekind dan mengajukan masalah berikut:
\begin{quote}
Dapatkah suatu permukaan (misalnya suatu persegi beserta batasnya)
dikorelasikan satu-satu dengan suatu garis (misalnya suatu garis lurus
beserta titik-titik ujungnya), sehingga setiap titik permukaan berpadanan
dengan suatu titik garis dan, sebaliknya, setiap titik garis berpadanan
dengan suatu titik permukaan?

Saat ini masih tampak bagi saya bahwa jawaban atas pertanyaan ini sangat
sulit---meskipun di sini pun orang begitu terdorong untuk menjawab
\emph{tidak}, sehingga orang ingin menganggap pembuktiannya hampir tidak
perlu. [Dikutip dalam \citealt{Gouvea2011}]
\end{quote}
Namun, pada 1877, Cantor membuktikan bahwa ia keliru. Sesungguhnya, suatu
garis dan suatu persegi mempunyai jumlah titik yang tepat sama. Pada
29 Juni 1877, ia menulis kepada Dedekind,
``\emph{je le vois, mais je ne le crois pas}''; yakni,
``Saya melihatnya, tetapi saya tidak memercayainya''. Dalam
``sejarah yang lazim diterima'' mengenai matematika, ucapan ini sering
dianggap menunjukkan betapa \emph{sungguh sulit dipercayanya} hasil-hasil
baru tersebut bagi para matematikawan pada masa itu. (Korespondensinya
disajikan dalam \citet{Gouvea2011}, dan kita kembali membahasnya dalam
\olref[his][set][mythology]{sec}. Bukti Cantor diuraikan secara ringkas
dalam \olref[his][set][cantorplane]{sec}.)

Terinspirasi oleh hasil Cantor, Peano mulai mempertimbangkan apakah mungkin
memetakan suatu garis \emph{secara kontinu} ke seluruh bidang. Ini akan
menjadi suatu \emph{kurva pengisi ruang}. Pada \citeyear{Peano1890}, Peano
mengonstruksi tepat kurva seperti itu. Hasil ini benar-benar berlawanan
dengan intuisi: Euclid mendefinisikan garis sebagai ``panjang tanpa
lebar'' (Buku I, Definisi 2), tetapi Peano menunjukkan bahwa dengan
melengkungkan suatu garis secara tepat, panjangnya dapat diubah menjadi
lebar. Pada \citeyear{Hilbert1891}, Hilbert menguraikan suatu kurva pengisi
ruang yang agak lebih intuitif, beserta beberapa gambar yang
mengilustrasikannya. Kurva itu dikonstruksi secara berurutan; berikut enam
tahap pertama konstruksinya:
\begin{center}
\begin{tikzpicture}
\begin{scope}[xshift=20pt, yshift=20pt]
	\draw [oldiagcolorC, l-system={Hilbert curve, step=40pt, angle=90, axiom=L, order=1}] lindenmayer system;
\end{scope}
\begin{scope}[xshift=110pt, yshift=10pt]
	\draw [oldiagcolorC, l-system={Hilbert curve, step=20pt, angle=90, axiom=L, order=2}] lindenmayer system;
\end{scope}
\begin{scope}[xshift=205pt, yshift=5pt]
	\draw [oldiagcolorC, l-system={Hilbert curve, step=10pt, angle=90, axiom=L, order=3}] lindenmayer system;
\end{scope}
\begin{scope}[xshift=2.5pt, yshift=-97.5pt]
	\draw [oldiagcolorC, l-system={Hilbert curve, step=5pt, angle=90, axiom=L, order=4}] lindenmayer system;
\end{scope}
\begin{scope}[xshift=101.25pt, yshift=-98.75pt]
	\draw [oldiagcolorC, l-system={Hilbert curve, step=2.5pt, angle=90, axiom=L, order=5}] lindenmayer system;
\end{scope}
\begin{scope}[xshift=200.625pt, yshift=-99.375pt]
	\draw [oldiagcolorC, l-system={Hilbert curve, step=1.25pt, angle=90, axiom=L, order=6}] lindenmayer system;
\end{scope}
\draw[gray] (0pt,0pt) rectangle (80pt, 80pt);
\draw[gray] (100pt,0pt) rectangle (180pt, 80pt);
\draw[gray] (200pt,0pt) rectangle (280pt, 80pt);
\draw[gray] (0pt,-100pt) rectangle (80pt, -20pt);
\draw[gray] (100pt,-100pt) rectangle (180pt, -20pt);
\draw[gray] (200pt,-100pt) rectangle (280pt, -20pt);
\end{tikzpicture}
\end{center}
Pada limit---suatu pengertian yang saat itu telah menerima definisi
ketat---seluruh persegi terisi penuh dengan !!{colorC}. Sebagai catatan,
kurva Hilbert kontinu di mana-mana tetapi tidak dapat didiferensialkan di
mana pun; secara intuitif karena, pada limit tak hingga, fungsi itu mengubah
arah secara mendadak pada setiap saat. (Kita akan menguraikan konstruksi
Hilbert secara lebih terperinci dalam
\olref[his][set][hilbertcurve]{sec}.)

Entah baik atau buruk, konstruksi geometris ``patologis'' ini diperlakukan
sebagai alasan untuk meragukan penggunaan intuisi geometris. Konstruksi itu
menjadi sesuatu yang mendekati \emph{propaganda} bagi cara baru mengerjakan
matematika, yang kelak memuncak pada teori himpunan. Dalam pembentukan mitos
tentang bidang ini pada masa berikutnya, sering diulangi bahwa hasil-hasil
tersebut sepenuhnya ketat sekaligus sepenuhnya mengejutkan. Karena itu,
hasil-hasil tersebut melayani dua tujuan: sebagai peringatan agar tidak
bergantung pada intuisi geometris, dan sebagai demonstrasi kesuburan cara
berpikir baru.

\end{document}
